\documentclass[11pt, letterpaper]{article}
\usepackage{amsmath, amssymb}
\usepackage{hyperref}
\usepackage{booktabs}
\usepackage{enumitem}
\usepackage{xcolor}
\usepackage{geometry}
\usepackage{fancyhdr}
\usepackage{titlesec}
\usepackage{tcolorbox}
\usepackage{graphicx}
\usepackage{caption}
\usepackage{subcaption}
\usepackage{parskip}
\geometry{margin=1in}

\definecolor{courseblue}{HTML}{002452}
\definecolor{coursegold}{HTML}{FABD0F}
\definecolor{coursered}{HTML}{B90E31}
\definecolor{lightblue}{HTML}{E8EEF7}
\definecolor{lightgold}{HTML}{FEF8E3}

\hypersetup{colorlinks=true,linkcolor=courseblue,urlcolor=coursered}
\tcbuselibrary{skins}
\newtcolorbox{guidancebox}{enhanced,colback=lightblue,colframe=courseblue,
  sharp corners,leftrule=4pt,left=6pt,right=6pt,top=4pt,bottom=4pt}

\pagestyle{fancy}\fancyhf{}
\fancyhead[L]{\small\textcolor{courseblue}{\textbf{CISC 867 --- Deep Learning | Fall 2026}}}
\fancyhead[R]{\small\textcolor{courseblue}{Final Report}}
\fancyfoot[C]{\small\thepage}
\renewcommand{\headrulewidth}{0.4pt}
\titleformat{\section}{\large\bfseries\color{courseblue}}{}{0em}{}[\titlerule]
\titleformat{\subsection}{\normalsize\bfseries\color{courseblue}}{}{0em}{}

\begin{document}

{\centering
  {\Large\textbf{CISC 867: Deep Learning --- Final Report}}\\[0.5em]
  {\large\textit{[PROJECT ID AND TITLE]}}\\[0.4em]
  \textbf{Group Members:}\\
  {[Student 1 Name]} (Project Lead) \quad
  {[Student 2 Name]} (Research Lead) \quad
  {[Student 3 Name]} (Engineering Lead)\\[0.3em]
  \textbf{GitHub:} \url{[link]} \quad \textbf{Date:} \today\\[0.4em]
  \rule{\linewidth}{0.4pt}
}

\begin{abstract}
\begin{guidancebox}
\textbf{Guidance.} 200--250 words. State: (1) the problem, (2) your method,
(3) your main result compared to baselines, and (4) your novel contribution.
\end{guidancebox}
[WRITE YOUR ABSTRACT HERE. Include your best headline number.]
\end{abstract}

\section{Introduction}

\begin{guidancebox}
\textbf{Guidance.} 1/2--1 page. Motivate the problem, situate it in the
broader DL literature, and end with a clear bullet list of your contributions.
The last paragraph \emph{must} read: ``In summary, our contributions are:''
followed by a \texttt{itemize} list.
\end{guidancebox}

%[MOTIVATION AND CONTEXT.]

In summary, our contributions are:
\begin{itemize}[noitemsep]
  \item Contribution 1, e.g., ``We reproduce X on datasets A and B, achieving Y.''
  \item Contribution 2, e.g., ``We propose Z, which improves performance by W\%.''
  \item Contribution 3, e.g., ``We identify failure mode F and propose mitigation M.''
\end{itemize}

\section{Related Work}

\begin{guidancebox}
\textbf{Guidance.} 1/2--1 page. Synthesise the related work into 2--3 paragraphs
by theme (e.g., ``CNN-based approaches'', ``Transformer-based approaches'').
Explicitly distinguish your work from each group. Cite $\geq$7 papers.
\end{guidancebox}

%[RELATED WORK BY THEME, WITH EXPLICIT COMPARISONS TO YOUR WORK.]

\section{Method}

\begin{guidancebox}
\textbf{Guidance.} 1--1.5 pages. Describe your approach with enough precision
to reproduce it. Include: (1) a model diagram or pseudocode, (2) mathematical
notation for the key components, (3) training objective, (4) any modifications
to an existing method. Use the notation you use in your equations.
\end{guidancebox}

\subsection{Problem Formulation}

[DEFINE THE TASK FORMALLY. E.g., ``Given input $x \in \mathcal{X}$, we seek a
mapping $f_\theta$ such that $\hat{y} = f_\theta(x)$ minimises...'']

\subsection{Model Architecture}

[DESCRIBE YOUR ARCHITECTURE. Reference a figure if you include one.]

\subsection{Training Objective}

[STATE YOUR LOSS FUNCTION WITH EQUATION. E.g.,
\begin{equation}
  \mathcal{L} = \lambda_1 \mathcal{L}_{\mathrm{task}} + \lambda_2 \mathcal{L}_{\mathrm{aux}}
  \label{eq:loss}
\end{equation}
where\ldots]

\subsection{Proposed Modification}

[DESCRIBE WHAT YOU CHANGED OR ADDED RELATIVE TO THE BASELINE.
WHY DID YOU EXPECT THIS TO HELP? WHAT IS THE HYPOTHESIS?]

\section{Experiments}

\begin{guidancebox}
\textbf{Guidance.} This section carries the most weight. Structure it as:
setup $\to$ main results $\to$ ablation $\to$ qualitative analysis $\to$
failure analysis. All quantitative results must be reported as mean $\pm$ std
over $\geq$3 independent seeds with different random initialisations.
\end{guidancebox}

\subsection{Experimental Setup}

\begin{tabular}{ll}
\toprule
\textbf{Setting} & \textbf{Value} \\
\midrule
Dataset(s) & [e.g., CIFAR-10: 50k train / 10k test] \\
Metrics & [e.g., Top-1 Accuracy, ECE, AUROC] \\
Optimiser & [e.g., AdamW, $\beta_1=0.9$, $\beta_2=0.999$] \\
Learning rate & [e.g., $5\times10^{-4}$ with cosine decay] \\
Batch size & [e.g., 128] \\
Training duration & [e.g., 100 epochs] \\
Seeds & [e.g., 42, 123, 999] \\
Hardware & [e.g., CPU: Intel i7-12700H; or Google Colab T4 GPU] \\
\bottomrule
\end{tabular}

\subsection{Main Results}

\begin{guidancebox}
\textbf{Guidance.} Provide one main results table comparing your method to
$\geq$2 published baselines. Each cell: mean $\pm$ std from 3 runs. Bold the best.
Write 2--3 paragraphs of analysis below the table.
\end{guidancebox}

\begin{table}[h]
\centering
\begin{tabular}{lcc}
\toprule
\textbf{Method} & \textbf{Metric 1} $\uparrow$ & \textbf{Metric 2} $\downarrow$ \\
\midrule
Baseline 1 (Pub.) & [Published value] & [Published value] \\
Baseline 2 (Pub.) & [Published value] & [Published value] \\
Our reproduction & [Value] $\pm$ [Std] & [Value] $\pm$ [Std] \\
\textbf{Ours (proposed)} & \textbf{[Value] $\pm$ [Std]} & \textbf{[Value] $\pm$ [Std]} \\
\bottomrule
\end{tabular}
\caption{Main results on [DATASET], [SPLIT]. All reported results are mean $\pm$ std over 3 seeds.
$\uparrow$ = higher is better; $\downarrow$ = lower is better.}
\label{tab:main}
\end{table}

[ANALYSIS: Interpret Table~\ref{tab:main}. What do the numbers show?
Is the improvement statistically meaningful? What is the effect size?]

\subsection{Ablation Study}

\begin{guidancebox}
\textbf{Guidance.} Ablate each component of your proposed method individually.
Start from the full model and remove one thing at a time, or start from the
baseline and add one thing at a time. Each ablation should have a clear
hypothesis and result.
\end{guidancebox}

\begin{table}[h]
\centering
\begin{tabular}{lccc}
\toprule
\textbf{Configuration} & \textbf{Component A} & \textbf{Component B} & \textbf{Metric} \\
\midrule
Full model (ours) & \checkmark & \checkmark & [Value] $\pm$ [Std] \\
w/o Component A & $\times$ & \checkmark & [Value] $\pm$ [Std] \\
w/o Component B & \checkmark & $\times$ & [Value] $\pm$ [Std] \\
Baseline & $\times$ & $\times$ & [Value] $\pm$ [Std] \\
\bottomrule
\end{tabular}
\caption{Ablation study on [DATASET]. Each row is mean $\pm$ std over 3 seeds.}
\label{tab:ablation}
\end{table}

[ANALYSIS: What does Table~\ref{tab:ablation} show about which components
contribute most? Are the contributions additive? Independent?]

\subsection{Qualitative Analysis}

[INSERT FIGURE(S) SHOWING REPRESENTATIVE PREDICTIONS. For image tasks: a grid
of input/output/ground-truth comparisons. For text tasks: example outputs.
For classification: a confusion matrix or t-SNE plot.]

\begin{figure}[h]
  \centering
  % \includegraphics[width=\linewidth]{figures/qual.png}
  \caption{[CAPTION: describe what each panel shows and what it demonstrates.]}
  \label{fig:qual}
\end{figure}

\subsection{Failure Analysis}

\begin{guidancebox}
\textbf{Guidance.} Identify $\geq$3 categories of failure cases. For each:
show an example (or describe one), explain \emph{why} the model fails,
and suggest a potential mitigation. This section is highly valued by reviewers.
\end{guidancebox}

\paragraph{Failure type 1:} [NAME] --- [DESCRIBE WHEN AND WHY THE MODEL FAILS.
HOW WOULD YOU MITIGATE THIS?]

\paragraph{Failure type 2:} [NAME] --- [DESCRIBE.]

\paragraph{Failure type 3:} [NAME] --- [DESCRIBE.]

\section{Discussion}

\begin{guidancebox}
\textbf{Guidance.} 1/2 page. Interpret your results in the context of the
literature. Were your hypotheses confirmed? What did you learn?
What would you do differently with more time or compute?
\end{guidancebox}

[INTERPRETATION AND DISCUSSION. CONNECT BACK TO YOUR RESEARCH QUESTION.]

\subsection{Limitations}

[HONESTLY LIST THE LIMITATIONS OF YOUR WORK: compute constraints, dataset
size, hyperparameter sensitivity, evaluation scope, etc.]

\section{Conclusion}

[CONCISE SUMMARY: what you set out to do, what you found, and one concrete
future direction for follow-up work.]

\section{AI Tool Disclosure}

\begin{guidancebox}
\textbf{Required for all submissions.} Per course policy, you must declare
all AI tools used and how you used them.
\end{guidancebox}

[EXAMPLE: ``We used Claude (Anthropic, claude.ai) for help debugging our
training loop (Sec.~4) and for writing feedback on Sec.~5. All analysis,
interpretation, and experimental results are our own. We verified all
AI-generated claims against primary sources.'']

\section{Contribution Declaration}

\begin{tabular}{lcc}
\toprule
\textbf{Member} & \textbf{Contributions} & \textbf{\%} \\
\midrule
{[Student 1]} & [e.g., Method design, report Secs.~3--4, experiments] & \% \\
{[Student 2]} & [e.g., Literature review, ablations, figures] & \% \\
{[Student 3]} & [e.g., Data pipeline, code, README, reproducibility] & \% \\
\bottomrule
\end{tabular}

\section{References}

\begin{enumerate}[label={[\arabic*]},leftmargin=*]
  \item [Author(s), Year. ``Title.'' \textit{Venue}. arXiv:XXXX.XXXXX]
  \item [\ldots]
\end{enumerate}

\appendix

\section{Appendix: Additional Results}

\begin{guidancebox}
\textbf{Guidance.} Place supplementary tables, figures, and extended ablations
here. The appendix does not count toward the 6--8 page main body limit.
\end{guidancebox}

[ADDITIONAL RESULTS, EXTENDED ABLATIONS, IMPLEMENTATION DETAILS.]

\section{Appendix: Reproducibility Checklist}

Confirm each item before submission:

\begin{itemize}[noitemsep]
  \item[$\square$] \texttt{README.md} explains how to reproduce all main results with one command.
  \item[$\square$] \texttt{requirements.txt} or \texttt{environment.yml} lists all dependencies with versions.
  \item[$\square$] All random seeds are fixed and documented.
  \item[$\square$] All results are mean $\pm$ std over $\geq$3 independent runs.
  \item[$\square$] GitHub repository is public and the link is in the report header.
  \item[$\square$] No hard-coded absolute paths in any script.
  \item[$\square$] All datasets are downloaded programmatically or clearly documented.
  \item[$\square$] Model checkpoints or inference instructions are included (if model is too large, document how to retrain from scratch).
\end{itemize}

\end{document}
